\begin{otherlanguage}{ngerman}

\chapter*{Zusammenfassung}

NextGenPB ist ein Finite-Elemente-basierter Löser für die lineare Poisson-Boltzmann-Gleichung mit Fokus auf Recheneffizienz, der auf das Hochleistungsrechnen ausgerichtet ist, GPUs jedoch bislang nicht nutzt, um die Leistung noch weiter zu steigern.
In dieser Arbeit habe ich verschiedene Techniken eingesetzt, um ihn mit Nvidias CUDA-API zu beschleunigen.
Der Schwerpunkt lag auf der Lösung des linearen Gleichungssystems sowie auf der Energieberechnung.
Einige zusätzliche hostseitige Optimierungen verringerten die Gesamtlaufzeit weiter und verbesserten zugleich die Speichereffizienz.

Für das lineare Gleichungssystem habe ich Nvidias Löser-Bibliothek AMGX mit großem Erfolg eingesetzt und dabei Beschleunigungen von bis zu 10x erreicht.
Ich habe zwei Konfigurationen erstellt, eine mit Fokus auf Speichereffizienz und eine mit Fokus auf Geschwindigkeit; beide übertrafen LIS bei Weitem.
Dabei zeigte sich, dass AMGX über mehrere GPUs innerhalb eines Knotens gut skaliert, nicht jedoch über mehrere Knoten hinweg, wo der Kommunikationsaufwand die Gewinne zunichtemachte oder sogar umkehrte.

Für die Energieberechnung habe ich eine naive Implementierung und eine schnelle Multipolmethode entwickelt.
Die naive Implementierung brachte für die Energieberechnung des Moleküls 1VSZ eine enorme Beschleunigung von 24,6x auf einer RTX 3080 und von 321x auf einer A100.
Darüber hinaus bot die schnelle Multipolmethode für compute-gebundene Probleme eine noch größere Beschleunigung und steigerte den Faktor auf der RTX 3080 auf 117x, wobei die ursprüngliche Genauigkeit von NextGenPB erhalten blieb.
Zwar verliert die schnelle Multipolmethode ihren Vorteil bei nicht compute-gebundenen Problemen, dies betrifft jedoch nur Probleme, die ohnehin sehr wenig Zeit benötigen.

Diese Optimierungen machten es möglich, das gesamte H1N1-Kapsid auf einem einzelnen Knoten mit 4 H200-GPUs in unter einer Stunde zu berechnen.
Sie ermöglichen es zudem, mittelgroße Probleme mit weniger Clusterknoten zu berechnen.

\end{otherlanguage}
